\title{Direct Preference Optimization: Your Language Model is Secretly a Reward Model}

\begin{abstract}
Large-scale unsupervised language models (LMs) acquire extensive knowledge and certain reasoning capabilities, yet exercising fine-grained control over their outputs remains challenging because their training lacks direct supervision. To address this, current approaches gather human assessments of the relative quality of generated outputs and then adapt the LM to reflect these preferences, most commonly through reinforcement learning from human feedback (RLHF). Nonetheless, RLHF involves a multi-stage and frequently unstable process: it first requires training a reward model to capture human judgments, and then employs reinforcement learning to optimize the original LM for this learned reward, all while constraining deviation from the initial model parameters. This work proposes an alternative reward model formulation for RLHF that permits direct closed-form derivation of the corresponding optimal policy, effectively reducing the standard RLHF task to minimization of a straightforward classification loss. The proposed approach, termed \textit{Direct Preference Optimization} (DPO), is robust, effective, and computationally efficient, removing the necessity for generation sampling during fine-tuning and greatly reducing the burden of hyperparameter selection. Experimental results indicate that DPO aligns LMs with human preferences at least as effectively as prior techniques. Significantly, DPO-based fine-tuning surpasses PPO-driven RLHF in managing generational sentiment and delivers equal or improved performance in tasks like summarization and single-turn dialogue, all with a far simpler implementation and training process.
\end{abstract}

\section{Introduction}
Large-scale unsupervised language models (LMs) trained on immense corpora have demonstrated unexpected strengths~\citep{chowdhery2022palm, brown2020language, touvron2023llama, bubeck2023sparks}. Yet, these systems are trained using text sourced from humans with highly diverse motivations, values, and levels of expertise. Naturally, not every goal or skill found in the training data should be emulated; for instance, while an AI-based coding assistant must \textit{recognize} frequent programming errors to address them, it should, during generation, prefer the (possibly infrequent) segments of exemplary programming proficiency embedded within its dataset. Likewise, for factual accuracy, we might want a language model to be \textit{cognizant} of widely held misconceptions, but it is undesirable for it to reproduce such inaccuracies with the same frequency as their prevalence in the human population. Thus, it is imperative to distinguish and select the model's \emph{targeted actions and answers} from its broad \textit{range of knowledge and skills} if we wish to construct AI systems that are safe, reliable, and controllable~\citep{ouyang2022training}. Although reinforcement learning (RL) is the prevalent approach for aligning LMs with human preferences, we demonstrate that the standard RL-based learning objective can be matched precisely by employing a straightforward binary cross-entropy loss, thus streamlining the overall process of learning from preferences.

In general, contemporary techniques introduce preferred behaviors to a language model through curated collections of human preference data that exemplify attributes humans consider safe and beneficial. This stage of preference learning follows an initial phase of massive unsupervised pre-training over large-scale text datasets. While the most direct form of preference alignment consists of supervised fine-tuning with high-quality human-generated exemplars, the prevailing and more successful set of methods centers on reinforcement learning from human (or AI) feedback, commonly known as RLHF/RLAIF~\citep{christiano2017deep, bai2022constitutional}. RLHF protocols entail fitting a reward model to pairwise preference judgments and then employing RL to steer the LM policy to generate outputs with high reward, all while limiting deviation from the initial pre-trained model. Despite the impressive linguistic and programming prowess of RLHF-fine-tuned models, this approach entails significantly greater complexity than pure supervised learning, requiring the training of several language models and repeated online sampling during training, which leads to substantially higher computational requirements.

This work introduces a method to train language models to follow human preferences directly, avoiding the need for explicit reward modeling or reinforcement learning techniques. We present \textit{Direct Preference Optimization} (DPO), an approach that implicitly maximizes the same objective as standard RLHF frameworks—that is, reward maximization subject to a KL-divergence penalty—yet remains easy to implement and efficient to train. Conceptually, DPO modifies the model by increasing the log-odds of favorable responses compared to unfavorable ones, utilizing dynamically adjusted, example-specific importance weights. These weights serve to mitigate issues of model collapse seen with a simple probability ratio approach. As with traditional methods, DPO depends on a formal model of preferences, such as the Bradley-Terry model \cite{bradley1952rankanalysis}, to assess the concordance between a reward function and observed preference data. Nevertheless, instead of first learning a reward function via a preference loss and subsequently optimizing the policy to maximize this learned reward, DPO leverages a change of variables to re-express the preference loss directly in terms of the policy. This allows DPO, given human-labeled preference data for model responses, to optimize the policy using a straightforward binary cross entropy loss—yielding an optimal policy for a latent reward function matched to the preference observations.

The principal contribution of this study is Direct Preference Optimization (DPO), a streamlined, reinforcement learning–free method for optimizing language models from preference data. Empirical results demonstrate that DPO is at least as effective as current leading approaches, such as PPO-based RLHF, in preference-centric tasks including sentiment transfer, summarization, and conversational dialogue, even when applied to models as large as 6B parameters.

\section{Related Work}

Self-supervised language models have demonstrated the capacity to handle certain tasks in a zero-shot manner \citep{radford2019language}, or with only a handful of examples as prompts \citep{gpt3,megatron,chowdhery2022palm}. Nevertheless, their effectiveness on downstream applications and alignment with user intent are markedly enhanced when they are further trained on corpora containing explicit instructions paired with human-authored completions \citep{mishra-etal-2022-cross,sanh2022multitask,chung2022scaling,thoppilan2022lamda}. This process, known as ‘instruction-tuning’, permits large language models to generalize beyond the set of seen instructions and generally improves their overall utility \citep{chung2022scaling}. While instruction tuning has yielded considerable progress, acquiring \textit{relative} human assessments of response quality tends to be more scalable and practical than collecting expert-level demonstrations. Accordingly, recent studies have adopted the practice of fine-tuning large language models on datasets that reflect human preferences, leading to improved model capabilities in areas such as translation \citep{kreutzer-etal-2018-reliability}, summarization \citep{stiennon2022learning,ziegler2020finetuning}, creative writing \citep{ziegler2020finetuning}, and following instructions \citep{ouyang2022training,ramamurthy2023is}. Typically, these approaches involve first fitting a neural reward model to agree with human preference data using a preference modeling framework, like the Bradley-Terry model \citep{bradley1952rankanalysis}. Subsequently, reinforcement learning methods—most often REINFORCE \citep{williams1992reinforce}, proximal policy optimization (PPO; \cite{schulman2017proximal}), or their variants \citep{ramamurthy2023is}—are employed to adapt the language model so as to maximize the learned reward. Related research takes advantage of instruction-following LLMs fine-tuned with human feedback to create synthetic preference datasets targeting aspects like safety or harmlessness \citep{bai2022constitutional}, relying solely on high-level text rubrics for guidance rather than direct, granular supervision. This synergy merges progress from two distinct domains: one focusing on applying reinforcement learning to language model training for a variety of goals~\citep{Ranzato2015SequenceLT,paulus2018a,wu2018learning}, and another addressing general strategies for learning from human preference judgments \citep{christiano2017deep,kupcsik2018learning}. Even though learning from comparative human feedback is highly attractive, the reinforcement learning fine-tuning of large language models continues to pose significant practical difficulties. In this work, we propose a theoretically grounded method to optimize for relative human preferences that circumvents the need for reinforcement learning.

Beyond linguistic applications, the study of policy learning from preferences has been explored within both bandit frameworks and reinforcement learning environments, leading to a variety of methodological advances. One such direction is contextual bandit learning that utilizes preferences or ordinal rankings over actions in place of explicit reward signals, commonly referred to as the contextual dueling bandit (CDB) paradigm \citep{yue2012karmed, dudik2015contextual}. In the theoretical analysis of CDBs, given the absence of scalar reward information, the criterion of optimality is supplanted by the concept of the \textit{von Neumann winner}: a policy that is guaranteed to have an expected victory rate of at least 50\% when compared with any other policy \citep{dudik2015contextual}. A salient distinction is that, in CDB settings, preference information is typically acquired online, whereas human preference learning settings more often involve training from a static offline collection of preference-labeled action pairs \citep{yan2022human}. Similarly, in \textit{preference-based RL} (PbRL), learning proceeds from binary preferences that are induced by an unknown ‘scoring’ or utility function, bypassing explicit reward feedback \citep{BusaFekete2014, ruiz2023dueling}. The suite of existing PbRL algorithms includes those capable of leveraging off-policy preference datasets, but these generally require first estimating an underlying scoring function (that is, the reward model) before proceeding to policy optimization \citep{jain2013learning, BusaFekete2014, christiano2017deep, sadigh2017active, kupcsik2018learning}. In contrast, our proposed approach unifies policy learning and preference satisfaction into a single optimization step by directly updating the policy in accordance with observed preferences.

\section{Preliminaries}\label{section:prelims}

We begin by summarizing the RLHF pipeline described by \citeauthor{ziegler2020finetuning}, with additional reference to later work (\citep{stiennon2022learning, bai2022training, ouyang2022training}). This standard pipeline encompasses three key stages: (1) supervised fine-tuning (SFT), (2) preference-based sampling and reward modeling, and (3) policy optimization via reinforcement learning.

\textbf{SFT}: The RLHF process is typically initiated by taking a pre-trained language model and further training it on curated, task-specific datasets through supervised learning. This yields the supervised fine-tuned model $\pi^\text{SFT}$, which serves as the basis for subsequent preference modeling (e.g., for dialogue or summarization tasks).

\textbf{Reward Modelling Phase}: During the next phase, the SFT model is prompted with input prompts $x$, and generates paired candidate outputs $(y_1, y_2) \sim \pi^\text{SFT}(y \mid x)$. These output pairs are evaluated by human annotators, who indicate a preference between the alternatives. Formally, this is expressed as $y_w\succ y_l \mid x$, where $y_w$ and $y_l$ denote the selected (preferred) and rejected (less preferred) response among $(y_1, y_2)$, respectively. These observed preferences are assumed to stem from a hidden reward function $r^*(y, x)$, to which the system does not have direct access. Several methodologies exist for modeling such preference data; the Bradley-Terry (BT) model \cite{bradley1952rankanalysis} is commonly used, though more general approaches such as the Plackett-Luce model \citep{plackett1975analysis, luce2012individual} can also be employed when richer ranked outputs are available. Under the BT framework, the distribution of human preferences $p^*$ is represented as follows:

Given access to a fixed dataset of pairwise comparisons $\mathcal{D}=\bigl\{x^{(i)}, y_w^{(i)}, y_l^{(i)}\bigr\}_{i=1}^N$ drawn from $p^*$, a reward function $r_{\phi}(x, y)$ can be modeled, with parameters estimated by maximizing likelihood. By casting this setup as a binary classification task, we obtain the following negative log-likelihood loss:

\begin{equation}\label{eq:reward_model}
    \mathcal{L}_R(r_{\phi}, \mathcal{D}) = -\mathbb{E}_{(x, y_w, y_l)\sim \mathcal{D}}\bigl[\log \sigma(r_{\phi}(x, y_w)- r_{\phi}(x, y_l))\bigr]
\end{equation}

where $\sigma$ refers to the logistic sigmoid. In applications involving LMs, the architecture $r_{\phi}(x, y)$ typically inherits its weights from an SFT model $\pi^\text{SFT}(y \mid x)$, with a single linear head appended to the terminal transformer layer to yield a scalar reward estimate \cite{ziegler2020finetuning}. To decrease the variance of the resulting reward estimates, previous work enforces normalization such that  $\mathbb{E}_{x,y\sim \mathcal{D}}\left[r_\phi(x, y)\right] = 0$ for all $x$.

\textbf{RL Fine-Tuning Phase}: In this phase, feedback is provided to the language model through the learned reward. In line with existing approaches~\citep{jaques2017sequence, jaques2020human}, the objective takes the following form:

\begin{equation}\label{eq:RL}
\max_{\pi_{\theta}}  \mathbb{E}_{x\sim \mathcal{D}, y\sim \pi_{\theta}(y \mid x)}\bigl[r_{\phi}(x, y)\bigr] - \beta\mathbb{D}_{\textrm{KL}}\bigl[\pi_{\theta}(y\mid x)\mid \mid \pi_\text{ref}(y\mid x)\bigr],
\end{equation}

Here, $\beta$ serves as a regularization coefficient that constrains divergence from the reference policy $\pi_\text{ref}$, which is commonly taken to be the original SFT model $\pi^\text{SFT}$. 
Typically, the language model policy $\pi_\theta$ is initialized to match $\pi^\text{SFT}$. Incorporating this KL penalty is crucial to ensure that the learned policy does not stray excessively from the data distribution on which the reward model remains reliable, while also promoting diversity in generated outputs and mitigating collapse onto a few high-reward responses. Because generating text involves sampling from a discrete distribution, this objective cannot be optimized via direct differentiation, necessitating reinforcement learning techniques. The prevailing solution \citep{ziegler2020finetuning, stiennon2022learning, bai2022training, ouyang2022training} is to define the modified reward ${r(x, y) = r_{\phi}(x, y) -\beta (\log \pi_{\theta}(y\mid x) - \log \pi_\text{ref}(y\mid x))}$, and optimize using PPO \cite{schulman2017proximal}.

\section{Direct Preference Optimization}\label{sec:DPO}

Acknowledging the complexity of applying RL algorithms at the scale required for language model fine-tuning, we seek an alternative framework for optimizing policies directly based on observed preferences. Distinct from earlier RLHF techniques—which first fit an explicit reward model and subsequently optimize it with RL—our methodology employs a particular reward parameterization permitting direct extraction of the optimal policy in closed form, sidestepping the necessity for reinforcement learning altogether. As detailed in the following sections, the central idea involves exploiting a known analytical correspondence between reward functions and their associated optimal policies, thereby converting a loss defined over rewards into one over policies. This change-of-variable approach removes the need for a standalone reward network while preserving fidelity to established human preference modeling paradigms like the Bradley-Terry model. Consequently, the policy model is simultaneously responsible for both language generation and implicitly encoding the reward.

\textbf{DPO Objective Derivation.} Our starting point aligns with previous work, employing the RL objective from Eq.~\ref{eq:RL} under an arbitrary reward function $r$. As established in prior literature~\citep{peters2007reinforcement, peng2019advantage, korbak2022reinforcement, go2023aligning}, the maximization of a reward function subject to a KL-divergence constraint yields the optimal policy:

\begin{equation}\label{eq:op_policy}
    \pi_r(y\mid x) = \frac{1}{Z(x)}\pi_\text{ref}(y\mid x)\exp\left(\frac{1}{\beta}r(x, y)\right),
\end{equation}

where $Z(x) =\sum_{y}\pi_\text{ref}(y\mid x)\exp\left(\frac{1}{\beta}r(x, y)\right)$ denotes the normalization constant. A full derivation can be found in Appendix \ref{app:derivation1}. Despite the possibility of leveraging the MLE-estimated reward function $r_{\phi}$ as a surrogate for the true reward $r^*$, the computation of the partition function $Z(x)$ remains challenging and computationally burdensome~\citep{korbak2022reinforcement, go2023aligning}, hindering the practical use of Eq.~\ref{eq:op_policy}. To address this, we can rewrite Eq.~\ref{eq:op_policy}, isolating $r(x, y)$ in terms of $\pi_r$, $\pi_\text{ref}$, and $Z(\cdot)$. Specifically, applying the logarithm to both sides and rearranging gives:

\begin{equation}\label{eq:main_eq}
    r(x,y) =\beta \log \frac{\pi_r(y\mid x)}{\pi_\text{ref}(y\mid x)} + \beta \log Z(x).
\end{equation}

This reparameterization is valid for both the ground-truth reward $r^*$ and its associated optimal policy $\pi^*$. Notably, the Bradley-Terry framework utilizes the difference in rewards for two responses, specifically, ${p^*(y_1 \succ y_2 \mid x) = \sigma(r^*(x, y_1) - r^*(x, y_2))}$. Substituting the form of $r^*(x, y)$ from Eq.~\ref{eq:main_eq} into the preference model, Eq.~\ref{eq:bradley-terry}, causes the partition function to drop out. As a result, the probability of human preferences can be expressed solely through the optimal policy $\pi^*$ and the reference policy $\pi_\text{ref}$, independent of any reward parameterization. Thus, under the Bradley-Terry model, the RLHF-optimal policy $\pi^*$ adheres to the following relationship:

\begin{equation}\label{eq:objective}
    p^*(y_1\succ y_2 \mid x)=\frac{1}{1 + \exp\left(\beta \log \frac{\pi^*(y_2\mid x)}{\pi_\text{ref}(y_2\mid x)} - \beta \log \frac{\pi^*(y_1\mid x)}{\pi_\text{ref}(y_1\mid x)}\right)}
\end{equation}

A step-by-step derivation is provided in Appendix~\ref{app:derivation2}. While this expression leverages the Bradley-Terry preference model, similar developments can be performed for other models, such as those of Plackett-Luce~\citep{plackett1975analysis, luce2012individual}, with derivations located in Appendix~\ref{app:plackett_luce_models}.

This reformulation now allows us to directly connect the likelihood of human preference outcomes to the policy, bypassing explicit reward modeling. Consequently, a maximum likelihood objective can be constructed for a parameterized policy $\pi_\theta$. In analogy to conventional reward modeling (cf. Eq.~\ref{eq:reward_model}), the policy optimization objective is:

\begin{equation}\label{eq:optimum_model}
    \mathcal{L}_\text{DPO}(\pi_{\theta}; \pi_\text{ref}) = -\mathbb{E}_{(x, y_w, y_l)\sim \mathcal{D}}\left[\log \sigma \left(\beta \log \frac{\pi_{\theta}(y_w\mid x)}{\pi_\text{ref}(y_w\mid x)} - \

In this framework, we approximate an implicit reward function via an alternative parametrization such that the optimal policy coincides with $\pi_\theta$. Furthermore, our approach amounts to learning a reparameterized version of the Bradley-Terry model, which imparts advantageous theoretical guarantees—including consistency—provided that the preference data distribution satisfies certain conditions \cite{bong2022generalized}. Further theoretical implications and comparisons with related methodologies are provided in Section~\ref{sec:theory}.

\textbf{What mechanism underlies the DPO update?} To deepen our mechanistic comprehension of DPO, let us examine the gradient of its loss, $\mathcal{L}_\text{DPO}$. The gradient with respect to the parameter vector $\theta$ can be expressed as follows:

\begin{multline*}\label{eq:gradient}
    \nabla_\theta \mathcal{L}_\text{DPO}(\pi_\theta;\pi_\text{ref}) = \\ -\beta\mathbb{E}_{(x, y_w, y_l) \sim \mathcal{D}} \bigg[\underbrace{\sigma(\hat{r}_\theta(x, y_l) - \hat{r}_\theta (x, y_w))}_\text{higher weight when reward estimate is wrong}\bigg[\underbrace{\nabla_\theta\log \pi(y_w \mid x)}_\text{increase likelihood of $y_w$} - \underbrace{\nabla_\theta\log\pi(y_l \mid x)}_\text{decrease likelihood of $y_l$}\bigg]\bigg],
\end{multline*}

Here, $\hat{r}_\theta(x, y) = \beta \log \frac{\pi_\theta(y \mid x)}{\pi_\text{ref}(y \mid x)}$ serves as the reward implicitly determined by the current policy $\pi_\theta$ in relation to the reference $\pi_\text{ref}$ (see Section~\ref{sec:theory} for further elaboration). Conceptually, the gradient update of $\mathcal{L}_\text{DPO}$ acts to boost the likelihood of the favored completion $y_w$, while suppressing that of the less favored $y_l$. Notably, the magnitude of these adjustments depends on the extent to which the implicit reward model $\hat{r}_\theta$ mistakenly prefers $y_l$ over $y_w$, modulated by $\beta$—thus reflecting the severity of the ranking error as well as the intensity of the KL regularization. Empirical evidence from our studies highlights the criticality of this importance weighting: omitting the weighting factor (i.e., employing a simpler version of the update) can result in deleterious model behaviors, such as language degeneration (see Appendix Table~\ref{tab:unlikelihood_generations}).

\textbf{DPO overview.} The standard DPO process comprises the following steps: 1) For each prompt $x$, generate response pairs $y_1, y_2 \sim \pi_\text{ref}(\cdot \mid x)$, then use human evaluations to annotate these pairs and create a dataset of preference judgments $\mathcal{D} = \{x^{(i)}, y_w^{(i)}, y_l^{(i)}\}_{i=1}^N$; 2) train the language model $\pi_\theta$ by minimizing $\mathcal{L}_\text{DPO}$ with respect to the chosen $\pi_\text{ref}$, dataset $\mathcal{D}$, and the hyperparameter $\beta$.

Practically, existing publicly released preference datasets are preferred over creating new samples and collecting additional human annotations. These datasets are often generated using $\pi^\text{SFT}$, so when such a policy is available, we set $\pi_\text{ref} = \pi^\text{SFT}$. If no such supervised fine-tuned model $\pi^\text{SFT}$ exists, we estimate $\pi_\text{ref}$ by maximizing the likelihood over preferred completions ${(x, y_w)}$: specifically, ${\pi_\text{ref} = \argmax_{\pi}\mathbb{E}_{x, y_w \sim \mathcal{D}}\left[\log \pi(y_w \mid x)\right]}$. This initialization step reduces mismatch between the true (but unobserved) reference distribution and the practical $\pi_\text{ref}$ used by DPO. Appendix~\ref{app:implementation} contains further implementation specifics and hyperparameter information.

\section{Theoretical Analysis of DPO} This section provides a deeper theoretical exploration of the DPO approach, elucidates its interpretative framework, and discusses its theoretical strengths—particularly in comparison to actor-critic methods in RLHF like PPO~\cite{schulman2017proximal}.

\label{sec:theory}

\subsection{Your Language Model Is Secretly a Reward Model} DPO enables simultaneous policy and reward optimization using a single maximum likelihood framework, thereby obviating the need for separate reward modeling and reinforcement learning steps. The core optimization goal, presented in Eq. \ref{eq:main_eq}, mirrors a Bradley-Terry model, with rewards defined as $r^*(x, y) = \beta \log\frac{\pi^*_\theta(y \mid x)}{\pi_\text{ref}(y \mid x)}$. Optimization over the parameterized model $\pi_{\theta}$ can thus be viewed, through a suitable change of variables, as optimizing a reward model as in Eq. \ref{eq:reward_model}. In what follows, we systematically develop the underlying theory of this parameterization, demonstrating that it fully preserves the expressiveness of the reward model class and permits exact retrieval of the optimal policy. Our analysis begins with the formal definition of equivalence for reward functions.

\begin{definition}
We define two reward functions $r(x, y)$ and $r'(x, y)$ as equivalent if and only if ${r(x, y)-r'(x, y) = f(x)}$ for some function $f$.     
\end{definition}

One can immediately verify that this defines an equivalence relation, which organizes the collection of reward functions into distinct equivalence classes. We now present two lemmas:

\begin{lemma}\label{lemma:same_prefrence} Within the Plackett-Luce framework, including the Bradley-Terry model, any pair of reward functions belonging to the same equivalence class will induce identical preference distributions.
\end{lemma}

\begin{lemma}\label{lemma:same_policy}
    For the constrained RL setting, two reward functions from a given equivalence class generate the same optimal policy.
\end{lemma}

The arguments supporting these lemmas are direct and are included in Appendix \ref{app:lemma1}. The first lemma is a recognized issue concerning the identifiability within the Plackett-Luce family \cite{plackett1975analysis}, often referred to as under-specification. This under-specification typically necessitates introducing further identifiability assumptions to provide guarantees for maximum likelihood estimation, as described in Eq. \ref{eq:reward_model} \cite{bong2022generalized}. The second lemma asserts that all reward functions within a particular class induce the identical optimal policy. Therefore, in our analysis, it suffices to recover any representative reward function from the appropriate equivalence class. We formalize this idea in the following theorem, with its proof provided in Appendix~\ref{app:thm1}:

\begin{theorem}\label{thm:main}
    Given mild conditions, every reward equivalence class compatible with the Plackett-Luce (and the Bradley-Terry as a special case) models can be captured by the parameterization ${r(x, y) = \beta \log \frac{\pi(y\mid x)}{\pi_\text{ref}(y\mid x)}}$, where $\pi(y\mid x)$ is some model and $\pi_\text{ref}(y \mid x)$ is a fixed reference model.
\end{theorem}

\begin{sproof}
    Let $r(x, y)$ denote a reward function which yields an associated optimal policy $\pi_r(y \mid x)$, as defined in Eq. \ref{eq:op_policy}. We aim to show that there exists a member of the equivalence class of $r$ expressible in the aforementioned reparameterized form. Consider the projection
\begin{equation}
    f(r; \pi_\text{ref}, \beta)(x, y) = r(x, y) - \beta\log\sum_{y}\pi_\text{ref}(y\mid x)\exp\left(\frac{1}{\beta}r(x, y)\right)
\end{equation}
This operator $f$ effectively standardizes the reward by subtracting the log-partition function (which depends only on $x$) computed with respect to $\pi_r$. Since the normalization is solely a function of the context $x$, $f(r; \pi_\text{ref}, \beta)(x, y)$ remains within the same equivalence class as $r(x, y)$. Substituting $r$ with the right-hand side of Eq.~\ref{eq:main_eq} (valid for any reward function), it follows that $f(r; \pi_\text{ref}, \beta)(x, y) = \beta \log \frac{\pi_r(y\mid x)}{\pi_\text{ref}(y\mid x)}$. Thus, $f$ yields a reward function of the desired form that belongs to the equivalence class of $r$, confirming that this reparameterization does not compromise the generality of the reward model.
\end{sproof}

Theorem~\ref{thm:main} can also be interpreted as uniquely determining the reward function chosen by the DPO reparameterization within each equivalence class. Specifically, it identifies the reward function that fulfills:

\begin{equation}\label{eq:lag_p}
     \sum_{y}\underbrace{\pi_\text{ref}(y\mid x)\exp\left(\frac{1}{\beta}r(x, y)\right)}_{=\pi(y\mid x)\text{, using Thm.~\ref{thm:main} reparam.}} = 1,
\end{equation}

meaning that $\pi(y\mid x)$ constitutes a legitimate probability distribution (non-negative probabilities adding up to one). Observing Eq.~\ref{eq:op_policy}, it becomes evident that Eq.~\ref{eq:lag_p} characterizes the partition function associated with the optimal policy generated by the reward function $r(x, y)$. The essential idea underlying the DPO method is the imposition of specific constraints on the otherwise underdetermined Plackett-Luce family (notably, the Bradley-Terry model among them) of preference models. This restriction preserves the expressive capacity regarding representable reward models, yet allows for an explicit, analytically tractable expression of the optimal policy in Eq.~\ref{eq:op_policy} across all possible prompts $x$.

\subsection{Instability of Actor-Critic Algorithms} The framework we have established also facilitates analysis of instabilities observed in conventional actor-critic methods within RLHF, such as PPO. Adhering to the RLHF protocol, our focus is on the reinforcement learning fine-tuning phase described in Section \ref{section:prelims}. There exists a connection here to the control as inference paradigm \cite{levine2018reinforcement}, particularly for the constrained RL formulation presented in \ref{eq:RL}. Letting $\pi_{\theta}(y\mid x)$ denote the policy parameterization, the optimization seeks to minimize $\mathbb{D}_{\text{KL}}[\pi_{\theta}(y|x) \mid \mid \pi^*(y\mid x)]$, where $\pi^*$—the optimal policy—arises from Eq.~\ref{eq:optimum_model} via the reward function $r_{\phi}(y, x)$. Basic manipulations yield the following optimization target:

\begin{equation}\label{eq:AC}
    \max_{\pi_{\theta}}\mathbb{E}_{\pi_{\theta}(y\mid x)}\bigg[\underbrace{r_{\phi}(x, y) -\beta\log\sum_{y}\pi_\text{ref}(y\mid x)\exp\left(\frac{1}{\beta}r_{\phi}(x, y)\right)}_{f(r_{\phi}, \pi_\text{ref}, \beta)} - \underbrace{\beta\log\frac{\pi_{\theta}(y\mid x)}{\pi_\text{ref}(y\mid x)}}_{\text{KL}}\bigg]
\end{equation}

The objective function under consideration aligns with those optimized in earlier studies 
\citep{ziegler2020finetuning, stiennon2022learning, bai2022training, ouyang2022training}, employing a DPO-equivalent reward structure for the $r_{\phi}$ reward class. Within this framework, the normalization component of $f(r_{\phi}, \pi_\text{ref}, \beta)$ can be seen as the soft value function associated with the reference policy $\pi_\text{ref}$. Although the inclusion or exclusion of this term does not influence the solution that optimizes the objective, omitting it may increase the variance of the policy gradient estimator, leading to greater instability during training. One solution is to learn the normalization term via an estimated value function, though such an approach may introduce its own optimization difficulties. Historically, other approaches have employed a human completion baseline for reward normalization, which is effectively a single-sample Monte Carlo approximation of the normalizing factor. By contrast, the DPO reparameterization produces a reward function inherently free from the need for such baselines.

\section{Experiments}
This section presents a systematic assessment of DPO's capacity to train policies directly using preference data. Initially, we explore a controlled text-generation domain, investigating DPO’s efficiency in balancing the maximization of rewards with the minimization of KL-divergence relative to a reference policy. This comparison is made against established preference optimization methods such as PPO. Subsequently, we analyze DPO’s effectiveness on larger models and more complex RLHF challenges, including summarization and dialogue tasks. Our empirical results indicate that, with minimal hyperparameter adjustment, DPO frequently matches or outperforms strong baseline methods—such as RLHF utilizing PPO—as well as surpassing approaches based on selecting the best of $N$ sampled trajectories evaluated under a learned reward. Prior to discussing these empirical findings, we outline the experimental protocols; further experimental specifics can be found in Appendix~\ref{app:exp_details}.

\textbf{Tasks.} We conduct experiments across three distinct open-domain text generation scenarios. In every case, the algorithms are trained to optimize a policy based on a dataset of pairwise preferences, $\mathcal{D}=\bigl\{x^{(i)}, y_w^{(i)}, y_l^{(i)}\bigr\}_{i=1}^N$. In the \textbf{controlled sentiment generation} task, $x$ represents the beginning of a movie review drawn from the IMDb dataset \cite{maas-EtAl:2011:ACL-HLT2011}, and the model must complete it with a continuation $y$ expressing positive sentiment. To facilitate rigorous evaluation, preference pairs over model outputs are \textit{automatically} produced using a pre-trained sentiment classifier, ensuring that $p(\text{positive}\mid x,y_w) > p(\text{positive}\mid x,y_l)$. For the SFT baseline, GPT-2-large is fine-tuned on the training portion of the IMDB dataset until performance stabilizes (details provided in App~\ref{app:sentiment_details}). For the \textbf{summarization} task, $x$ is a Reddit forum submission and the model is required to output a summary $y$ that captures the key points. In alignment with previous research, we employ the Reddit TL;DR summarization dataset \citep{volske-etal-2017-tl} and supplement it with human preference annotations from \citeauthor{stiennon2022learning}. Our SFT baseline here involves a model refined on human-produced Reddit summaries\footnote{\url{https://huggingface.co/CarperAI/openai_summarize_tldr_sft}} using the TRLX \citep{leandro_von_werra_2023_7790115} RLHF toolkit. The human preference pairs originate from \citeauthor{stiennon2022learning} and were annotated on outputs of a different, though similarly fine-tuned, SFT model. Lastly, the \textbf{single-turn dialogue} task uses $x$ as a human prompt, ranging from technical queries to requests for personal advice. The model must generate an appropriate, useful reply $y$; for this, we draw from the Anthropic Helpful and Harmless dialogue dataset \citep{bai2022training}, which comprises 170,000 exchanges between people and an AI assistant. Each dialogue concludes with two machine-generated responses and a human judgment indicating the preferred option. In this case, as a pre-existing SFT model does not exist, we build an SFT model by further tuning a standard language model solely on the responses preferred by humans.

\textbf{Evaluation.} We adopt two main strategies for evaluation across our experiments. To assess how effectively each algorithm optimizes the constrained reward maximization objective, we utilize a controlled sentiment generation scenario. Here, each method is evaluated by examining its Pareto frontier in terms of both attained reward and KL-divergence from the base policy, leveraging the availability of an explicit ground-truth reward signal provided by a sentiment classifier. In practical, real-world scenarios where this ground-truth reward is inaccessible, our evaluation instead centers on the \textit{win rate} metric relative to a baseline, employing GPT-4 as a stand-in for human evaluators to judge the quality of summaries and the helpfulness of responses in summarization and single-turn dialogue tasks, respectively. For the summarization task, the baseline is the test set’s reference summary; for the dialogue setting, it is the response indicated as preferred within the test set. While prior work indicates that large language models may serve as more reliable automatic evaluators than standard metrics \citep{Chen2023ExploringTU}, we supplement this claim by conducting a human evaluation, detailed in Sec.~\ref{sec:human-judgments}, to substantiate our reliance on GPT-4. Our results show a strong correlation between GPT-4 and human ratings, with the degree of human alignment with GPT-4 matching or surpassing inter-annotator consistency.

\textbf{Methods.} Beyond DPO, we include a variety of established strategies for training language models to reflect human preferences in our evaluation. Specifically, for the summarization task, we use zero-shot prompts with \textbf{GPT-J} \citep{gpt-j}, and for the dialogue task, 2-shot prompts with \textbf{Pythia-2.8B} \citep{biderman2023pythia}. We further evaluate the \textbf{SFT} model alongside \textbf{Preferred-FT}, which involves supervised fine-tuning on the selected response $y_w$ either generated by the SFT model (in both sentiment control and summarization) or by a generic language model (in the case of single-turn dialogue). Another approach is \textbf{Unlikelihood}~\citep{welleck2019neural}, which optimizes the model to increase the likelihood of $y_w$ while simultaneously decreasing the likelihood of $y_l$; here, an adjustable coefficient $\alpha\in[0,1]$ governs the strength of the unlikelihood objective. We also examine \textbf{PPO} \citep{schulman2017proximal}, where the reward signal is inferred from preference annotations, and \textbf{PPO-GT}, a form of oracle learning that uses the ground-truth reward available within the controlled sentiment setting. For the sentiment experiments, PPO-GT is implemented in two versions: one utilizes the standard implementation \cite{leandro_von_werra_2023_7790115}, while the other applies reward normalization and further hyperparameter tuning for optimal performance—adjustments also mirrored in ‘regular’ PPO runs with learned rewards. Finally, we assess a \textbf{Best of $N$} baseline, which samples $N$ outputs from the SFT (or Preferred-FT for dialogue) and chooses the top response as scored by a reward model trained on preference data. Although this method can yield strong results by separating reward model quality from PPO optimization, it is computationally expensive, as generating $N$ samples for every test prompt is not scalable even for moderate values of $N$.

\subsection{How well can DPO optimize the RLHF objective?}

The reward maximization objective constrained by KL divergence, which is standard in RLHF methods, seeks to maximize rewards without allowing the policy to diverge excessively from a reference. As such, evaluating and comparing algorithms requires considering both the level of reward attained and the corresponding KL divergence; obtaining marginally higher reward at the expense of significantly increased KL is not optimal. Figure~\ref{fig:frontier-tldr-main} presents the trade-off curves between reward and KL for various algorithms within the sentiment domain. For each algorithm, several training runs are performed, each using different settings for the policy's conservativeness hyperparameter (for PPO, target KL values of $\{3,6,9,12\}$; for DPO, $\beta$ values of $\{0.05,0.1,1,5\}$; for unlikelihood, $\alpha$ values of $\{0.05,0.1,0.5,1\}$; and for preferred-FT, different random seeds), summing to a total of 22 training runs. After every 100 training steps, up to convergence, each resulting policy is assessed using a suite of test prompts, calculating both the mean reward as given by the true reward function and the mean sequence-level KL\footnote{Computed as the aggregate of per-timestep KL-divergences.} with respect to the reference policy $\text{KL}\left(\pi\mid \mid \pi_\text{ref}\right)$. The experiments reveal that DPO establishes the most efficient frontier, simultaneously obtaining high reward and maintaining a low KL. This outcome stands out for several reasons. Notably, despite DPO and PPO sharing the same underlying objective, DPO consistently outperforms PPO in terms of the reward/KL efficiency; the trade-off curve for DPO strictly subsumes that for PPO. Furthermore, DPO manages to produce a superior frontier even in comparison to PPO-GT, where PPO has access to the actual reward signal.

\subsection{Can DPO scale to real preference datasets?}
\label{sec:dpo-real-datasets}
We proceed to assess the fine-tuning effectiveness of DPO on tasks involving summarization and single-turn dialogue. In the case of summarization, automatic evaluation metrics such as ROUGE often display weak alignment with human judgments~\citep{stiennon2022learning}, and earlier studies have demonstrated that fine-tuning language models via PPO on human preference data can yield superior summaries. To evaluate various approaches, we generate model completions from the test partition of the TL;DR summarization dataset and measure the mean win rate compared to reference completions within the test set. Completions are sampled for every method across temperatures ranging from 0.0 to 1.0, and the resulting win rates are illustrated in Figure~\ref{fig:frontier-tldr-main} (right). All methods, including DPO, PPO, and Preferred-FT, are fine-tuned from an identical GPT-J SFT base model\footnote{\url{https://huggingface.co/CarperAI/openai_summarize_tldr_sft}}. Our results indicate that DPO attains a win rate near 61\% at temperature 0.0, surpassing PPO’s peak performance of roughly 57\% at the same temperature. Additionally, DPO reaches a higher best-case win rate than the top $N$ baseline. It is important to mention that we did not perform significant tuning of DPO’s $\beta$ hyperparameter, suggesting these outcomes might not fully capture DPO's capacity. Furthermore, DPO demonstrates greater resilience to changes in sampling temperature compared to PPO, whose win rates can fall to those of the original GPT-J model when sampling at elevated temperatures. Preferred-FT yields only marginal gains over the baseline SFT model. A direct comparison between DPO and PPO based on human preference assessments is presented in Section~\ref{sec:human-judgments}, revealing that DPO outputs sampled at a temperature of 0.25 were favored in 58\% of cases relative to PPO outputs at the same temperature.

For single-turn dialogue tasks, we assess the various approaches using a portion of the Anthropic HH dataset’s test split \citep{bai2022training} that contains single rounds of interaction between human and assistant. Evaluation with GPT-4 leverages the preferred test completions as ground truth, calculating the win rate for each technique relative to these references. In the absence of a canonical SFT baseline for this domain, our procedure begins with a pre-trained Pythia-2.8B model; we apply Preferred-FT on selected completions to build a reference model such that its outputs reflect the data distribution, after which we train via DPO. Additionally, we benchmark against the top output among 128 completions produced by Preferred-FT (having observed that the Best of $N$ strategy ceases to improve beyond $N=128$ for this dataset, as detailed in Appendix Figure~\ref{fig:best-of-n}) as well as a 2-shot prompt version of the base Pythia-2.8B model. Across these setups, DPO consistently matches or exceeds the strongest-performing temperatures from all methods. We further evaluate an RLHF model fine-tuned with PPO on the Anthropic HH dataset \footnote{\url{https://huggingface.co/reciprocate/ppo_hh_pythia-6B}}, provided by a prominent community source \footnote{\url{https://github.com/CarperAI/trlx/tree/main/examples/hh}}, but find that neither the prompt nor the sampling temperature yields results that surpass the Pythia-2.8B baseline. Considering our TL;DR findings and the fact that both PPO and Best of 128 are geared toward maximizing the same reward function, we treat the Best of 128 baseline as an approximate indicator for PPO performance. In summary, DPO stands out as the only resource-efficient approach that delivers improvements over the Anthropic HH dataset's preferred completions and attains comparable or better outcomes than the computationally intensive Best of 128 baseline. As shown in Figure~\ref{fig:dialogue-main}, DPO also achieves its optimal performance in relatively few iterations.

\subsection{Generalization to a new input distribution}

To further investigate how PPO and DPO perform when facing shifts in input data, we assess the models trained on the Reddit TL;DR summarization task on an out-of-domain set: the test split of the CNN/DailyMail dataset \citep{nallapati-etal-2016-abstractive}, which contains news articles. We apply the optimal sampling temperatures identified on TL;DR (0 and 0.25) for evaluation. Table~\ref{tab:ood} shows the results. To maintain consistency with the original experimental protocol, GPT-4 win rates are calculated using a prompt similar to the one employed for Reddit TL;DR, substituting “forum post” with “news article.” On this unseen distribution, DPO maintains a substantial performance advantage over the PPO-trained policy. These findings provide preliminary support for the hypothesis that DPO-trained policies are capable of generalization on par with those trained via PPO, even though DPO does not utilize the supplemental unlabeled Reddit TL;DR prompts that PPO incorporates.

\subsection{Validating GPT-4 judgments with human judgments}
\label{sec:human-judgments}
To assess the trustworthiness of GPT-4’s evaluations, we conduct a human rating study using outputs from the TL;DR summarization task, along with two distinct GPT-4 prompts. The \textbf{GPT-4 (S)} (simple) prompt asks which summary better conveys the essential information in a post. In contrast, the \textbf{GPT-4 (C)} (concise) prompt also asks which summary is more succinct—a consideration motivated by our observation that GPT-4 tends to prefer lengthier, more repetitive outputs than humans when using the \textbf{GPT-4 (S)} prompt. Complete prompt texts are included in Appendix~\ref{app:prompts}. We design three pairwise comparisons using outputs from the best-performing method (DPO, temp. 0.25), the lowest-performing method (PPO, temp. 1.0), and a median-performing baseline (SFT, temp. 0.25), each evaluated against PPO with greedy decoding (its best temperature), to span a broad range of output quality. Our analysis shows that for both prompts, GPT-4’s preferences align with human preferences at a rate comparable to inter-human agreement, implying GPT-4 can serve as a valid surrogate for human judgment (for comparisons between DPO and PPO-1, multiple human assessments were gathered, though overall human rater availability was limited). Notably, the \textbf{GPT-4 (C)} prompt yields win rates most consistent with human evaluations; as such, this prompt is adopted for the principal analyses in Section~\ref{sec:dpo-real-datasets}. Further details on the human evaluation protocol, the interface used, and information on participating raters can be found in Appendix~\ref{app:human-study}.

\section{Discussion}
Preference-based learning serves as an effective and extensible method for building advanced, well-aligned language models. In this work, we proposed DPO, a straightforward approach for preference-based training of language models that eliminates the necessity for reinforcement learning. Instead of transforming the preference learning challenge into a classical RL framework simply to deploy existing RL algorithms, DPO constructs a direct correspondence between language model policies and reward structures, permitting language models to be optimized toward human preferences using a basic cross-entropy loss, sidestepping the complexities of reinforcement learning and retaining full generality. DPO matches or surpasses the performance of prevalent RLHF techniques such as those utilizing PPO, even with minimal hyperparameter adjustment, thereby substantially lowering the threshold for preference-based language model training.

\textbf{Limitations \& Future Work.} The outcomes of our study highlight numerous open avenues for further investigation. One question is the extent to which DPO-trained policies extrapolate to distributions beyond their training data in comparison with policies trained via explicit reward modeling. Preliminary experiments indicate that DPO's generalization capabilities are on par with PPO-trained policies, but a more thorough analysis is warranted. For instance, might leveraging DPO for self-labeling make unlabeled prompt datasets more useful? Another consideration concerns the phenomenon of reward over-optimization within the DPO framework and whether the minor decline in performance observed in Figure~\ref{fig:dialogue-main}-right can be attributed to this effect. Moreover, our empirical studies were limited to models up to 6B parameters; investigating the scalability of DPO for models on the scale of state-of-the-art, much larger systems remains an intriguing direction for future inquiry. In terms of evaluation methodology, we observed that GPT-4's computed win rates can vary depending on the prompt design; exploring robust techniques for extracting reliable assessments from automated evaluators is a topic for further research. Beyond these points, DPO could be applicable to a wide range of tasks, including the training of generative models across different data modalities, not just language modeling from human feedback.

\end{document}